\documentclass[12pt, a4paper]{article}

% --- Packages ---
\usepackage[margin=1in]{geometry}
\usepackage{amsmath, amssymb, amsthm}
\usepackage{setspace}
\usepackage{hyperref}
\usepackage{natbib}
\usepackage{booktabs}
\usepackage{enumitem}
\usepackage{titlesec}

% --- Theorem environments ---
\newtheorem{proposition}{Proposition}
\newtheorem{lemma}{Lemma}
\newtheorem{corollary}{Corollary}
\newtheorem{remark}{Remark}
\newtheorem{definition}{Definition}

% --- Formatting ---
\onehalfspacing
\setlength{\parskip}{0.6em}
\setlength{\parindent}{0pt}
\titleformat{\section}{\large\bfseries}{\thesection.}{0.5em}{}
\titleformat{\subsection}{\normalsize\bfseries}{\thesubsection.}{0.5em}{}

% --- Title ---
\title{\textbf{Referee Report: \\ ``Consumer Information and the Limits to Competition''} \\[6pt]
\large Mark Armstrong \& Jidong Zhou \\
\textit{American Economic Review}, 2022, 112(2), 534--577}
\author{Jordi Torres Vallverd\'{u}}
\date{\today}

\begin{document}

\maketitle
\thispagestyle{empty}

\section{Summary}

\subsection{Motivation}

In many markets, firms supply differentiated products and compete in prices, while consumers possess incomplete information about how well each product matches their preferences. A central question for platforms (e.g., Amazon, Booking) and regulators is how to design the information environment that consumers face.

The key tension is straight forward: more transparency improves product-consumer matching, but it also softens price competition by giving firms greater market power. Less transparency forces price competition --- products are perceived as near-perfect substitutes --- but consumers may end up purchasing a poorly matched product.

Armstrong and Zhou formalize this trade-off. They ask: within the class of all feasible signal structures ---provided by this external platform---, which one maximizes industry profit, and which one maximizes consumer surplus? Their answer reveals that the interests of firms and consumers are clearly opposed in how information should be structured, yet both sides agree that information should be symmetric across firms.

\subsection{Model Essentials}

Two symmetric firms ($i = 1,2$) supply differentiated products at zero cost. A risk-neutral consumer values product $i$ at $v_i \geq 0$; only her relative preference $x \equiv v_1 - v_2$ matters. Let $F$ denote the prior CDF of $x$, symmetric on $[-1,1]$, with density $f$. Two key parameters summarize the prior:
\begin{itemize}[nosep]
    \item $\mu \equiv E[v_i]$: expected value of a random product.
    \item $\delta \equiv E_F[\max\{x,0\}] = \int_{-1}^{0} F(x)\,dx$: expected match-quality gain from buying the preferred product over a random one.
\end{itemize}

Before buying, the consumer observes a private signal $s$ from a signal structure $\{\sigma(s|x), S\}$ chosen by a designer ---this is common knowledge. She updates to a posterior; since she is risk-neutral, only the posterior mean ($E[x \mid s]$) matters. Let $G$ denote the CDF of this posterior mean --- the posterior distribution. Firms observe $G$ but not $s$, and set prices simultaneously (Bertrand--Nash). The consumer buys from firm~1 if $E[x \mid s] > p_1 - p_2$.

The only constraint on $G$ is Bayesian consistency: $G$ must be a mean-preserving contraction (MPC) of $F$\footnote{Formally, this means: 
$$\int_{-1}^{x} G(\tilde{x})\,d\tilde{x} \;\leq\; \int_{-1}^{x} F(\tilde{x})\,d\tilde{x} \qquad \forall\, x \in [-1,1].$$ . Intuitively, a signal can not provide more information that what exists; it can only shift information across options.}.
Any such $G$ can be generated by some signal structure. So instead of optimizing over signals, the authors optimize directly over posteriors --- a key simplification.

Lemma 1 shows that a posterior $G$ can sustain symmetric equilibrium price $p$ iff $G$ lies between two bounds derived from each firm's no-deviation condition\footnote{Note this is only in pure strategies.}:
$$L_p(x) \;\leq\; G(x) \;\leq\; U_p(x)$$
Intuitively, the bounds rule out posteriors whose shape would make a unilateral price deviation profitable. A higher candidate price $p$ flattens both bounds, so that a larger set of posteriors can support it as an equilibrium.

%Consumer surplus under a symmetric posterior with price $p$ is $CS = \mu + \int_{-1}^{0} G(x)\,dx - p$. //no need to add this here., I think it adds nothing.


\subsection{Key Insight: Symmetric Signals Dominate}

Before solving the optimal design, the paper establishes a powerful restriction: for any asymmetric signal (one that favors a firm\footnote{We can think of this as a biased signal that tells consumers that one good is much better than the other.}), there exists a symmetric signal that yields higher profit for each firm and higher consumer surplus. Biased recommendations trigger aggressive undercutting by the disfavored firm, and this ends up hurting everyone. This means we can restrict attention to symmetric posteriors and symmetric prices throughout ---which is the second big simplification in the paper.

\subsection{Main Results}

\textbf{Firm-optimal signal.} The firm-optimal posterior sits on the lower bound of the feasibility corridor. The resulting density is U-shaped: consumers are pushed toward the extremes and few are near-indifferent, so neither firm can profitably undercut. Firms want the highest sustainable price, which means making the no-deviation constraint bind.

Two features of this outcome are key. First, there is \textbf{no mismatch}: the consumer buys her preferred product, and total welfare is maximized. Second, full information is not firm-optimal. Firms benefit from exaggerating perceived differentiation without changing which product the consumer prefers, which is what allows them to charge higher prices.

\textbf{Consumer-optimal signal.} The consumer-optimal posterior is less obvious. Consumers prefer low prices (which calls for pooling, no differentiation) but also prefer to buy their matched product (which calls for information). The paper shows that the price effect dominates: the optimal posterior brings the upper bound, with mass concentrated near $x = 0$. This creates many swing consumers and forces firms to compete fiercely, collapsing the price toward zero. The cost is \textbf{mismatch} --- some consumers buy the wrong product --- but the gain from lower prices outweighs it.

Importantly, the consumer-optimal signal is also not fully uninformative: some information can improve matches without relaxing price competition, as long as it keeps enough consumers near indifference.

% =============================================
% PART 2: CRITIQUE (1-2 pages)
% =============================================
\section{Critique}

I focus on three limitations which I find more substantial.\footnote{Two further points seem a bit second-order. First, the scalar reduction $x = v_1 - v_2$ is convenient but assumes full market coverage (no outside option), unit demand, and symmetric duopoly; once any of these fails, the level $v_1 + v_2$ matters and the one-dimensional framework no longer captures the full information-design problem. Second, the pure-strategy restriction is justified by the paper's 1.6\% bound of Proposition 4.}

\subsection{Single Designer and Commitment}

The paper thinks of the signal $G$ as what a platform can genuinely set: whether products are shown side by side, which attributes are displayed, how reviews are aggregated. Advertising and brand investment, which act on the prior, are not part of $G$ and are taken as given. Within this scope, two concerns remain.

\begin{itemize}[nosep]
    \item \textbf{Composite designer.} Even on-platform information is rarely the choice of a single agent. Third-party reviews, comparison sites, and consumer search all feed into the posterior alongside the platform's own display. The operative $G$ is then the outcome of a game among these sources, not a design choice, and the firm- and consumer-optimal benchmarks are not directly implementable.
    \item \textbf{Commitment.} The designer is assumed to commit to $G$ before prices are set. This matters asymmetrically for the two benchmarks. The consumer-optimal $G^{**}$ works precisely because it keeps consumers clustered near $x = 0$, which forces firms to compete aggressively. But a platform without commitment would want to reshape information ex post --- revealing a little more about each near-indifferent consumer lets firms soften competition (I discuss this in the extension).
\end{itemize}

\subsection{Symmetric and Exogenous Prior}

The prior over $(v_1, v_2)$ is shaped outside the model --- by advertising, brand investment, product positioning --- and the paper takes it as symmetric and fixed. Two concerns follow this simplifying assumption.

First, treating the prior as exogenous misses a margin of competition that interacts with signal design. Firms invest in branding to shift the prior in their favour, and those incentives depend on the signal they anticipate. A consumer-optimal $G$ that pools posteriors leaves little room to compete ex post, so firms have stronger reasons to invest ex ante in brand to secure a durable advantage. Signal design and prior formation act as substitutes for firms, and the paper's policy conclusions can unravel once firms respond off-platform.

Second, even holding the prior fixed, symmetry is doing real work. Lemma 2 (symmetric signals dominate) is derived from prior symmetry. With an asymmetric prior, posteriors that bias perception toward the ex ante stronger firm are no longer obviously dominated, and the reduction to symmetric $G$ and symmetric $p$ may break. Firms also disagree on what they want: the strong firm prefers signals that exaggerate its lead, the weak firm prefers to pool. The ``industry-profit'' benchmark that underlies the firm-optimal signal becomes ambiguous.

\subsection{Static Welfare and Dynamic Mismatch}

The consumer-optimal signal introduces mismatch by design: consumers near $x = 0$ buy a product they do not particularly prefer, and the paper treats this as a one-shot efficiency loss, $\delta - \int_{-1}^{0} G(x)\,dx$. In repeat settings, three further channels show up and each cuts against the static ranking.

\begin{itemize}[nosep]
    \item \textbf{Consumer learning.} A consumer who bought at $x \approx 0$ and regretted it will not be moved by the same signal next period. The more effective $G^{**}$ is today, the less it works tomorrow: the consumer-optimal design is dynamically inconsistent, and the platform would eventually have to reveal more to keep consumers engaged.
    \item \textbf{Public spillover.} Mismatched purchases generate returns and negative reviews, which feed into the prior that the \emph{next} cohort arrives with. Mismatch is not only a private loss but a negative externality on the information environment, and it brings back to the previous critique: the prior is not truly exogenous once we let reviews accumulate.
    \item \textbf{Softer competition through repeat purchase.} Firms that anticipate comeback buyers have weaker incentives to compete aggressively today. This partially undoes the price-pressure gain that made $G^{**}$ consumer-optimal in the first place.
\end{itemize}



\section{Extensions}

Armstrong and Zhou note in the conclusion that ``it would be valuable to embed this analysis within a framework in which the `information designer' is modeled explicitly as an economic agent''. I sketch two such extensions. I stay at the level of intuition: the goal is to identify what would change in the tradeoff between price competition and matching ---core tension of this paper---, and why this matters, rather than to re-derive the paper's results.

\subsection{Extension 1: Platform as Information Designer}

The natural candidate for an explicit designer is a \textbf{platform} --- Amazon, Expedia, an insurance comparator --- that controls how much product information consumers see and monetizes through fees. The question then becomes: which signal would a profit-maximizing platform choose, and how does its fee structure shape that choice?

Two fee instruments are the interesting cases: a per-unit (transaction) fee charged on every sale (e.g Booking charging transaction fees), and a lump-sum listing fee charged to firms upfront (e.g Booking charging a Hotel to be shown in the webpage).

\textbf{Lump-sum fees on firms.} Because a lump-sum fee does not enter marginal cost, it does not affect pricing. The platform can extract at most the industry profit $p(G)$, so it designs information to make that profit as large as possible. This is exactly the firm-optimal problem of the paper: the platform implements $G^*$, prices are high, there is no mismatch, and consumers keep the residual surplus $\mu + \delta - p^*$.

\textbf{Per-unit fees on consumers.} A per-unit fee is fully passed through, since demand is inelastic and allocation depends only on the price difference. The platform can then extract at most the consumer surplus generated by $G$, so it designs information to maximize $CS(G)$. This is the consumer-optimal problem: the platform implements $G^{**}$, prices are low, mismatch occurs, and consumers end up with zero net surplus --- the fee takes their gross gain away.

\textbf{The key point.} The platform's fee instrument selects a point on the Armstrong--Zhou frontier: charging firms makes the platform behave like a firm-side planner, charging consumers makes it behave like a consumer-side planner. Information design is not a primitive but an outcome of the business model. Two observations are worth flagging:

\begin{itemize}[nosep]
    \item Empirically, platform fees are usually firm-side (Amazon commissions, listing fees). Through this lens, this is also the case in which consumers are best off --- not because prices are low, but because the platform's incentive to maximize industry profit produces a no-mismatch outcome and leaves consumers the residual surplus.
    \item Adding a per-consumer fee is rare but not exotic (booking.com resort fees, Ticketmaster service fees, subscription platforms). The comparison shows that this seemingly ``consumer-facing'' instrument is precisely the one that erases consumer surplus in equilibrium.
\end{itemize}

\textbf{A subtle point on incentives.} One might worry that if firms anticipate the platform extracting $p(G^*)$ through the fee, they would just stop competing and set high prices. This does not actually happen in the static game, because the lump-sum fee is sunk once paid: it does not enter marginal pricing, so best responses and equilibrium prices are unchanged. What the fee \emph{does} change is participation --- firms must be left at least their outside option to log in to the platform --- and, in a richer model, investment incentives: if firms can invest in quality or advertising, a platform that extracts everything also destroys the incentive to invest, a textbook hold-up problem (``Amazon eats the margin''). The clean mapping between fee instrument and frontier point relies on ignoring these margins.

\textbf{Limitations.} The sketch assumes a monopoly platform with full commitment. Competing platforms would compete fees down and dilute the information-design incentive --- I suspect the duopoly case converges to something close to the non-strategic designer of the paper, while the monopoly platform pushes toward firm-best and captures the surplus through fees, but this is a natural next step rather than something I pin down here. Without commitment, the platform would adjust $G$ after observing prices and the timing game changes. A mixed-fee instrument lets the platform extract surplus from both sides, at which point the choice of $G$ is pinned down only by a no-mismatch condition.

\subsection{Extension 2: Asymmetric Prior (Vertical Differentiation)}

The paper's results rest heavily on prior symmetry. In many markets --- Coke vs.\ Pepsi, branded vs.\ generic, incumbent vs.\ entrant --- consumers arrive with an asymmetric prior: one product is perceived as better on average, shaped by habit, advertising, or past experience. I want to sketch what changes when the prior itself is tilted.

\textbf{Setup.} The cleanest way to formalize asymmetry is a mean shift: I keep the scalar $x = v_1 - v_2$ but now the prior $F$ on $x$ is centered at some $\mu_x > 0$ rather than $0$. Firm 1 is the ex ante stronger one. Bayesian consistency still imposes that any posterior $G$ is a mean-preserving contraction of $F$:
\[
\int_{-1}^{t} G(x)\,dx \;\le\; \int_{-1}^{t} F(x)\,dx \quad \text{for all } t,
\]
with equality at the endpoints. So the MPC toolkit of the paper still applies --- I just lose the convenient symmetry.

\textbf{Lemma 2 breaks.} Lemma 2 argues that industry profit is maximized by a symmetric $G$: given any asymmetric signal, the ``flipped'' version earns the same profit, and the average dominates. That swap only works because the prior itself is invariant under flipping. Once $\mu_x > 0$, the flipped signal no longer has the same prior and the argument collapses. Biased posteriors that amplify firm 1's lead are no longer automatically dominated.

\textbf{The ``firm-optimal'' benchmark becomes ambiguous.} In the symmetric case, Lemma 2 aligns the two firms: both prefer the same symmetric $G^*$, and ``industry profit'' is a meaningful object. Under an asymmetric prior this alignment breaks. The strong firm wants to exaggerate its lead: a posterior that pushes mass further to the right lets it raise its price while keeping most of the demand. The weak firm wants the opposite: posteriors that pool toward the indifference point, so that price competition remains alive and it retains positive demand. There is no longer a single ``firm-optimal'' signal --- the platform, or whoever designs information, has to pick \emph{whose} profit to maximize.

\textbf{The consumer-optimal twist.} I think this is the most interesting change. In the symmetric model, ``no information'' already produces fierce Bertrand-like competition because the prior sits exactly at the indifference point. With an asymmetric prior, ``no information'' gives firm 1 a structural advantage and prices are not competitive to begin with. So the consumer-optimal signal has to do something the symmetric case did not require: it has to actively \emph{correct} the asymmetry --- revealing more about the weak firm (to bring perceived valuations up) and pooling information about the strong firm (to bring it down), in order to cluster consumers around indifference and restore price competition.

Normatively this is nice: it gives a theoretical backing to policies that look odd through the paper's symmetric lens. Generic-drug labelling, mandatory comparative advertising, or regulators restricting branded-firm messaging can all be read as consumer-optimal information design under an asymmetric prior. The paper, by assuming symmetry, cannot say anything about these --- yet they are arguably the most common real-world interventions in the information environment.

\textbf{Link to the critique.} This also closes the loop with the advertising point from the critique. Much of what makes the prior asymmetric in the first place (brand investment, advertising, habit) happens off-platform. So the consumer-optimal on-platform signal is partly a response to off-platform distortions, and the two margins cannot really be studied separately. A parametric mean-shift prior would be the natural first pass; solving for $G^*$ and $G^{**}$ exactly is probably hard, but the qualitative ranking --- and whether the ``correction'' intuition survives --- should be tractable.

% =============================================
% REFERENCES
% =============================================
\section*{References}

\begin{itemize}[nosep, leftmargin=*]
    \item Armstrong, M. and J. Zhou (2022). ``Consumer Information and the Limits to Competition.'' \textit{American Economic Review}, 112(2), 534--577.
    \item Bergemann, D., B. Brooks, and S. Morris (2015). ``The Limits of Price Discrimination.'' \textit{American Economic Review}, 105(3), 921--957.
    \item Ivanov, M. (2013). ``Information Revelation in Competitive Markets.'' \textit{Economic Theory}, 52(1), 337--365.
    \item Johnson, J. and D. Myatt (2006). ``On the Simple Economics of Advertising, Marketing, and Product Design.'' \textit{American Economic Review}, 96(3), 756--784.
    \item Jullien, B. and A. Pavan (2019). ``Information Management and Pricing in Platform Markets.'' \textit{Review of Economic Studies}, 86(2), 893--952.
    \item Kamenica, E. and M. Gentzkow (2011). ``Bayesian Persuasion.'' \textit{American Economic Review}, 101(6), 2590--2615.
    \item Rochet, J.-C. and J. Tirole (2003). ``Platform Competition in Two-Sided Markets.'' \textit{Journal of the European Economic Association}, 1(4), 990--1029.
    \item Roesler, A.-K. and B. Szentes (2017). ``Buyer-Optimal Learning and Monopoly Pricing.'' \textit{American Economic Review}, 107(7), 2072--2080.
\end{itemize}
\end{document}
